\section*{結言}
この日は、AIモデルの能力競争が研究成果や開発基盤へ具体化する一方、エージェントが外部環境へ接続された際の統制が重要課題として浮上した。Claudeの理論物理計算やプラグイン基盤は専門領域とツール連携の拡大を示し、OpenAIの継続開示は訓練・評価環境であっても第三者サービスやユーザーデータへの影響を管理する必要性を示した。

生成メディアでは高解像度化とコード駆動制作が並行し、モデル比較も単純なベンチマークから実制作へ広がっている。政策・インフラ面でもリスク論と大規模計算基盤の拡張が同時に観測されており、能力向上、実運用、安全統制の三者が今後も同じ時間軸で進む流れが続きそうである。
